A Análise Exploratória de Dados é uma disciplina introduzida por John W. Turkey
na obra pioneira \textit{Explotarory Data Analysis} (1977) que visa a descoberta ou
identificação de aspectos ou padrões de interesse dentro de um conjunto de dados.
A AED surge como uma flexibilização dos métodos tradicionais da Estatística Descritiva.
Essa focava-se em resumir um conjunto de dados à determinadas medidas quantitativas,
empregando hipóteses menos flexíveis e sem a necessidade de generalizações,
geralmente com o objetivo de gerar relatórios.
Já a AED tenta empregar métodos mais gerais e flexíveis a fim de não somente resumir dados,
mas apresentá-los e representá-los de forma a destacar padrões e aspectos ocultos e gerar
novas hipóteses, extendendo as métricas descritivas com técnicas visuais.

Nesta parte, iremos abordar tanto as técnicas clássicas de resumo de dados como aquelas
que permitem a exploração de padrões e relações ocultas.
